\chapter*{Glossaire}
\addcontentsline{toc}{chapter}{Glossaire}
Les termes introduits dans les deux premiers chapitres sont réunis ici pour faciliter leur consultation.

\section*{Information et codage}
\begin{multicols}{2}\small\raggedright
\begin{rs}
\gloss{Pixel} : emplacement dans la grille d'une image, ou valeur qui y est stockée.
\gloss{Symbole} : élément de l'alphabet fini à coder.
\gloss{Mot de code} : suite de bits associée à un symbole.
\gloss{Code} : correspondance entre les symboles et leurs mots de code.
\gloss{Distribution empirique} : fréquences relatives observées des symboles.
\gloss{Histogramme} : représentation graphique de ces fréquences.
\gloss{Modèle de source} : loi de probabilité des symboles ou de leurs suites.
\gloss{Entropie} : information moyenne $H(V)$, positive ou nulle, mesurée en bits.
\gloss{Longueur moyenne du code} : $L(c)=\sum_vp_v\ell_v$, en bits par symbole.
\end{rs}
\end{multicols}

\section*{Traitement des images}
\begin{multicols}{2}\small\raggedright
\begin{rs}
\gloss{Bit et octet} : un chiffre binaire ; un groupe de huit bits.
\gloss{Canal} : l'un des tableaux de composantes d'une image multicanal.
\gloss{Quantification} : remplacement des valeurs par des indices d'un ensemble fini.
\gloss{Bruit} : fluctuations indésirables dues à l'acquisition ou au traitement.
\gloss{Moyenne} : somme des valeurs divisée par leur nombre.
\gloss{Médiane} : valeur centrale après classement d'un nombre impair de valeurs.
\gloss{Contour} : variation spatiale rapide des valeurs de l'image.
\gloss{Flou} : perte de transitions nettes et de détails fins.
\gloss{Contraste} : différences entre les régions claires et les régions sombres.
\gloss{Luminance} : grandeur photométrique définie à partir de la sensibilité visuelle.
\gloss{Matrice} : tableau rectangulaire $(a_{i,j})$.
\gloss{Résolution} : terme qui peut désigner les dimensions en pixels ou le pouvoir de séparation spatial ; ce sont deux notions distinctes.
\gloss{Fondu} : moyenne pondérée de deux images à positions de pixels fixes.
\end{rs}
\end{multicols}
